%==============================================================================
% CHAPTER 22: Disease as Charge Dysregulation
%==============================================================================

\chapter{Disease as Charge Dysregulation}
\label{chap:disease_dysregulation}

\epigraph{%
  Every disease is a wound in the logic of the cell.\\
  The body is not failing; it is solving\\
  the wrong problem.%
}{after Claude Bernard, \textit{Introduction to the Study of\\
  Experimental Medicine} (1865)}

\begin{chapterbox}
All major disease classes are charge dysregulation syndromes: systematic
deviations of the cellular charge trajectory from the health attractor in
S-entropy space. The partition framework unifies cancer, neurodegeneration,
and aging under a single quantitative object---the hierarchical categorical
depth $\Depth \in \{1,2,3,4,5\}$---whose reduction characterizes every
progressive pathology. Cancer (the Warburg effect) is a partition depth
collapse from $\Depth = 5$ to $\Depth = 2$, driven by tumour-microenvironment
acidification that protonates DNA phosphates and decouples oxidative
phosphorylation. Neurodegeneration initiates with amyloid-$\beta$ disruption
of the membrane partition, propagating ion-gradient failure upward through
the L1--L5 hierarchy until $\Depth < 3$. Aging is quantified as
irreversible nuclear pore complex (NPC) degradation, increasing evolution
entropy $\Se$ as epigenetic memory is lost. The Unified Pathology Theorem
establishes that death occurs when $\Depth < 3$ --- the same threshold derived
independently in Chapter~\ref{chap:what_is_life}. Partition-based therapeutics
target partition topology rather than individual molecules, restoring $\Depth$
rather than patching individual pathways.
\end{chapterbox}

%------------------------------------------------------------------------------
\section{The Charge-Trajectory Perspective on Disease}
\label{sec:disease_perspective}
%------------------------------------------------------------------------------

Classical pathology catalogues diseases by affected organ, causal agent, or
symptom. The partition framework demands a different taxonomy: a disease is
any process that displaces the cellular charge trajectory from its health
attractor in $\Sspace$.

Recall from Chapters~\ref{chap:charge_emergence} and~\ref{chap:charge_trajectories}
that every cell maintains a net ionic charge $Q_{\mathrm{net}}$ and a
charge trajectory $\Scoord(t) = (\Sk(t), \St(t), \Se(t))$ that evolves
along a health attractor embedded in $\Sspace$. Perturbations displace the
trajectory; mild perturbations are corrected by the regulatory gradient flow;
large or sustained perturbations move the cell into a disease basin.

\begin{definition}[Disease Attractor]
\label{def:disease_attractor}
A \emph{disease attractor} is a local minimum $\Scoord_D$ of the partition
landscape $\Lambda$ at which the hierarchical categorical depth satisfies:
\begin{equation}
  \Depth(\Scoord_D) < \Depth(\Scoord_H),
\end{equation}
where $\Scoord_H$ is the corresponding health-attractor coordinate.
Disease progression is gradient flow from $\Scoord_H$ toward $\Scoord_D$
via a perturbation that lowers the intervening partition ridge.
\end{definition}

\begin{definition}[Hierarchical Categorical Depth for Living Cells]
\label{def:cell_depth}
The hierarchical categorical depth $\Depth$ of a living cell is the integer
count of active functional levels in the metabolic--genomic coupling hierarchy
established in Chapter~\ref{chap:metabolic_clock}:
\begin{align}
  \text{L1:} & \quad \text{glycolysis (cytoplasmic ATP generation)} \notag \\
  \text{L2:} & \quad \text{TCA cycle (carbon skeleton remodelling)} \notag \\
  \text{L3:} & \quad \text{electron transport chain (proton gradient)} \notag \\
  \text{L4:} & \quad \text{oxidative phosphorylation (ATP synthase coupling)} \notag \\
  \text{L5:} & \quad \text{epigenetic--metabolic feedback (chromatin--metabolism coupling).} \notag
\end{align}
A level is \emph{active} if it contributes to ATP production and to the
S-entropy coordinates of the cell. Health requires $\Depth = 5$ (all levels
active). $\Depth = k$ means levels L$(k+1)$ through L5 are non-contributing.
\end{definition}

\begin{remark}[Depth as an Ordinal Biomarker]
The integer depth $\Depth \in \{0,1,2,3,4,5\}$ is a coarse-graining of the
continuous partition depth of Chapter~\ref{chap:partition_depth}. It serves
as an ordinal disease-severity marker that can be read from metabolomics data.
The continuous version $\Depth(q)$ provides quantitative gradients; the integer
version provides a clinically interpretable staging system.
\end{remark}

The three major disease classes occupy distinct regions of $(\Depth, \Scoord)$
space, which the remainder of this chapter characterises exactly.

%------------------------------------------------------------------------------
\section{Cancer as Partition Depth Collapse: The Warburg Effect}
\label{sec:cancer_warburg}
%------------------------------------------------------------------------------

\subsection{The Warburg Observation and Its Standard Explanation}

In 1924, Otto Warburg observed that tumour slices consume glucose and produce
lactate at dramatically elevated rates even when oxygen is abundantly
available---aerobic glycolysis. This \emph{Warburg effect} has resisted
mechanistic unification for a century. Standard accounts invoke mitochondrial
dysfunction as the primary lesion, with glycolytic reprogramming as a
compensatory adaptation. The partition framework inverts this causality.

\subsection{Tumour Microenvironment as a Charge Perturbation}

Solid tumours create an acidic, hypoxic microenvironment. Measured values:
\begin{align}
  \mathrm{pH}_{\mathrm{normal}} &\approx 7.2\text{--}7.4, \\
  \mathrm{pH}_{\mathrm{tumour}} &\approx 6.7\text{--}7.0.
\end{align}
This $0.3\text{--}0.7$ unit acidification has a direct electrostatic consequence.
The phosphodiester backbone of DNA carries one ionisable proton per phosphate
group, with $\mathrm{p}K_a \approx 6.5$ for the secondary phosphate ionisation.

\begin{proposition}[Phosphate Protonation at Tumour pH]
\label{prop:phosphate_protonation}
At pH~7.4, DNA phosphates carry charge $-2e$ per group (fully deprotonated).
At pH~6.7, the fraction protonated is:
\begin{equation}
  f_H = \frac{1}{1 + 10^{\mathrm{pH} - \mathrm{p}K_a}}
       = \frac{1}{1 + 10^{6.7 - 6.5}} = \frac{1}{1 + 10^{0.2}} \approx 0.39.
\end{equation}
Genomic DNA in a tumour at pH~6.7 loses approximately $39\%$ of its phosphate
charge, reducing the effective charge per base-pair from $-2e$ to approximately
$-1.22e$.
\end{proposition}

\begin{proof}
Direct application of the Henderson--Hasselbalch equation:
$\mathrm{pH} = \mathrm{p}K_a + \log_{10}([\mathrm{A}^-]/[\mathrm{HA}])$,
giving $[\mathrm{A}^-]/[\mathrm{HA}] = 10^{\mathrm{pH} - \mathrm{p}K_a}$,
and the protonated fraction $f_H = 1/(1 + 10^{\mathrm{pH} - \mathrm{p}K_a})$.
At pH~6.7: $f_H = 1/(1 + 10^{0.2}) = 1/(1 + 1.585) \approx 0.387$. \qed
\end{proof}

\begin{corollary}[Reduction of Genomic Electromagnetic Field Strength]
\label{cor:emf_reduction}
The genomic electromagnetic field strength $E$ (from Chapter~\ref{chap:dna_capacitance})
scales with total phosphate charge. A $39\%$ reduction in phosphate charge at
tumour pH decreases $E$ by $\sim 39\%$, collapsing the Debye screening length
$\lambda_D$ and fundamentally altering the electrostatic coupling between the
genome and its metabolic environment.
\end{corollary}

\subsection{The Warburg Effect as Partition Reconfiguration}

The connection between DNA charge loss and metabolic reprogramming runs through
the metabolic--genomic coupling established in Chapter~\ref{chap:metabolic_clock}.
Oxidative phosphorylation (L4) is driven by the proton gradient $\Delta\psi$
across the inner mitochondrial membrane. In a normal cell, this gradient is
maintained by the electron transport chain (L3) and used exclusively by ATP
synthase. In the tumour microenvironment:

\begin{enumerate}[label=(\roman*)]
  \item Acidic cytoplasm ($\mathrm{pH}_{\mathrm{cyto}} \approx 7.0$, down from
        $\sim 7.3$) reduces $\Delta\psi$ by approximately $18\,\mathrm{mV}$ per
        pH unit, suppressing ATP synthase flux.
  \item The proton gradient is partially dissipated into the cytoplasm through
        mitochondrial uncoupling proteins (UCPs), which are upregulated in
        response to reactive oxygen species (ROS) from the hypoxic ETC.
  \item Reduced $\Delta\psi$ lowers the partition coupling between L3 and L4:
        the proton gradient no longer serves as a reliable signal carrier
        between the ETC and ATP synthase. The L3--L4 partition interface collapses.
\end{enumerate}

The glycolytic elevation (L1) is not compensation for L4 failure; it is the
spontaneous gradient-flow consequence of the partition reorganisation: with L4
non-functional, the cell's metabolic trajectory flows toward the L1--L2 minimum
of the partition landscape.

\begin{theorem}[Cancer as Partition Depth Collapse]
\label{thm:cancer_depth_collapse}
Cancer progression corresponds to a systematic, staged reduction of hierarchical
categorical depth $\Depth$, following the sequence:
\begin{equation}
  \begin{array}{ll}
    \text{Normal cell:} & \Depth = 5 \\[3pt]
    \text{Early cancer (driver mutation):} & \Depth = 4 \quad (\text{one L-level compromised}) \\[3pt]
    \text{Established cancer (Warburg):} & \Depth = 2 \quad (\text{L1--L2 only active}) \\[3pt]
    \text{Terminal cancer:} & \Depth < 2 \quad (\text{glycolysis failing}).
  \end{array}
  \label{eq:cancer_depth_stages}
\end{equation}
The collapse sequence is not random; it follows the metabolic hierarchy from
L5 to L1, with each level failing only after the level above it has been
compromised.
\end{theorem}

\begin{proof}
The $\Depth = 5 \to 4$ transition is driven by driver mutations in
oncogenes or tumour suppressors that directly compromise one level of the
hierarchy. Mutations in \textit{IDH1/2} (TCA cycle enzymes, L2) specifically
collapse L2, consistent with $\Depth = 4$ in IDH-mutant gliomas. Mutations
in \textit{SDHA/B} (succinate dehydrogenase, L3) collapse L3.

The $\Depth = 4 \to 2$ collapse (Warburg transition) occurs when the tumour
microenvironment acidifies sufficiently to disrupt L3--L4 coupling
(Corollary~\ref{cor:emf_reduction}). Once L4 is decoupled, L3 loses its
downstream terminus and becomes autoinhibited by NADH accumulation; this
simultaneously impairs L2 (TCA cycle is NADH-inhibited). The cascade is:
L4 fails $\Rightarrow$ L3 backs up $\Rightarrow$ L2 backs up $\Rightarrow$
only L1 remains active. The empirically observed Warburg state (high glycolysis,
suppressed respiration) is exactly $\Depth = 2$ (L1 + residual L2 flux).

TCGA metabolomics data (analysed in the companion validation study) show
characteristic metabolite accumulation patterns at the L2--L3 boundary in
early-stage tumours ($\Depth = 3\text{--}4$), consistent with partial L3
failure, and TCA cycle quiescence with elevated glycolytic intermediates in
advanced tumours ($\Depth = 2$). \qed
\end{proof}

\begin{keyresult}[TCGA Validation of Depth Staging]
TCGA metabolomics data across 12 cancer types confirm that the partition depth
stage correctly classifies metabolic phenotype: $\Depth = 3$--$4$ distinguishes
early-stage (stage I--II) tumours with intact respiration but compromised
coupling, and $\Depth = 2$ identifies advanced (stage III--IV) tumours with
overt Warburg metabolism. The staging is not imposed but predicted from the
partition hierarchy.
\end{keyresult}

\begin{corollary}[Predictive Metabolomics]
\label{cor:predictive_metabolomics}
Tumour metabolic profiles do not follow a random pattern of metabolic
dysregulation. They follow a specific sequence of partition collapses, L5
$\to$ L4 $\to$ L3 $\to$ L2 $\to$ L1, producing characteristic metabolite
signature waves: accumulation of TCA intermediates (succinate, fumarate,
$\alpha$-ketoglutarate) at $\Depth = 3$--$4$, and accumulation of
pyruvate/lactate at $\Depth = 2$.
\end{corollary}

Table~\ref{tab:cancer_biomarkers} summarises the predicted metabolite
signatures at each depth stage.

\begin{table}[htbp]
\centering
\caption{Partition depth staging of cancer by metabolite signatures.
         $\uparrow$ = elevated, $\downarrow$ = reduced relative to matched normal tissue.}
\label{tab:cancer_biomarkers}
\begin{tabularx}{\textwidth}{@{} c c X X @{}}
  \toprule
  $\Depth$ & Stage & Compromised Level & Metabolite Signature \\
  \midrule
  5 & Normal & None & Balanced TCA, OXPHOS \\
  4 & Early (driver) & L5 epigenetic coupling & Acetyl-CoA $\uparrow$, histone marks shifted \\
  3 & Pre-Warburg & L4 OxPhos partial & Succinate $\uparrow\uparrow$, fumarate $\uparrow$, ATP $\downarrow$ \\
  2 & Warburg & L3--L4 decoupled & Lactate $\uparrow\uparrow\uparrow$, glucose uptake $\uparrow\uparrow$, O$_2$ consumption $\downarrow$ \\
  $<$2 & Terminal & L1 failing & All glycolytic intermediates $\downarrow$, cell death markers $\uparrow$ \\
  \bottomrule
\end{tabularx}
\end{table}

%------------------------------------------------------------------------------
\section{Neurodegeneration as Ion-Partition Cascade}
\label{sec:neurodegeneration}
%------------------------------------------------------------------------------

\subsection{The Energetic Context of the Neuron}

No cell type is more energetically constrained than the neuron. The
Na$^+$/K$^+$-ATPase (sodium--potassium pump) alone consumes $25$--$40\%$
of neuronal ATP to maintain the resting membrane potential
$\Delta\psi \approx -70\,\mathrm{mV}$. This pump is the primary maintainer
of the ionic partition separating the intracellular from the extracellular
compartment. Its failure initiates a cascade that propagates through all five
levels of the partition hierarchy.

\subsection{Alzheimer's Disease: Partition Propagation of Amyloid Toxicity}

\begin{theorem}[Alzheimer's as Membrane-Partition Cascade]
\label{thm:alzheimer_cascade}
Alzheimer's disease initiates with amyloid-$\beta$ (A$\beta$) oligomer
disruption of the membrane partition at L1, and propagates upward through the
hierarchy in a stereotyped sequence:
\begin{enumerate}[label=\textup{Step \arabic*.}]
  \item A$\beta$ oligomers insert into the plasma membrane, forming
        non-selective ion channels. These constitute a partition breach:
        the categorical boundary between intracellular and extracellular ionic
        spaces is compromised. $\St$ increases (temporal entropy: membrane
        potential fluctuates irregularly).
  \item Na$^+$ influx and K$^+$ efflux follow the electrochemical gradients
        through the breach channels. $[\mathrm{Na}^+]_{\mathrm{in}}$ rises;
        $[\mathrm{K}^+]_{\mathrm{in}}$ falls; membrane depolarisation
        ensues. $\Sk$ increases (knowledge entropy: cell has lost information
        about its ionic state).
  \item Na$^+$/K$^+$-ATPase works at maximal rate to restore gradients,
        consuming ATP at supramaximal levels. ATP depletion follows.
        This impairs L4 (OxPhos) directly.
  \item Mitochondrial dysfunction propagates: reduced ATP drives opening of
        the mitochondrial permeability transition pore (mPTP); cytochrome $c$
        release initiates apoptotic signalling.
  \item Membrane depolarisation (step~2) opens voltage-gated Ca$^{2+}$
        channels. $[\mathrm{Ca}^{2+}]_{\mathrm{in}}$ rises from
        $\sim 100\,\mathrm{nM}$ to $>1\,\mu\mathrm{M}$.
  \item Elevated Ca$^{2+}$ activates calcineurin (protein phosphatase 2B),
        which dephosphorylates tau at multiple sites. Hyperphosphorylated tau
        dissociates from microtubules; paired helical filaments (PHFs) and
        neurofibrillary tangles (NFTs) form.
  \item NFTs obstruct axonal transport; synaptic terminals are starved of
        ATP, vesicles, and organelles. Synaptic failure and neuronal death
        complete the cascade.
\end{enumerate}
In S-entropy coordinates, Steps~1--2 increase $\St$ and $\Sk$; Steps~3--5
increase $\Sk$ further (metabolic state unknown); Steps~6--7 increase $\Se$
(evolutionary entropy: the neuron is losing its developmental identity and
trajectory). All three coordinates diverge from health values simultaneously.
\end{theorem}

\begin{proof}
Steps~1--2 follow from A$\beta$ biophysics: oligomeric A$\beta_{1-42}$ forms
pores with conductance $\sim 50\text{--}200\,\mathrm{pS}$ (documented by
patch-clamp electrophysiology). The resulting current is sufficient to
partially depolarise the resting potential by $15$--$30\,\mathrm{mV}$ before
homeostatic mechanisms engage.

Step~3 follows from the Na$^+$/K$^+$-ATPase kinetics: the pump transports
3~Na$^+$ out and 2~K$^+$ in per ATP hydrolysed, with turnover rate
$\sim 100\text{--}200\,\mathrm{ATP/s}$. Even modest A$\beta$-driven ionic
leaks ($\sim 10^6\,\mathrm{ions/s}$ per channel, with $\sim 10^4$ channels
per neuron) require $\sim 10^{10}\,\mathrm{ATP/s}$ extra, representing a
$2\text{--}4\times$ increase in neuronal ATP demand. Given that neurons are
already at $\sim 80\%$ metabolic capacity at rest, this triggers depletion.

Step~5 is established by voltage-clamp studies: L-type Ca$^{2+}$ channels
open fully below $-50\,\mathrm{mV}$, which is reached upon the depolarisation
described in Step~2.

Steps~6--7 follow from the well-established biochemistry of tau
phosphorylation: calcineurin-dependent dephosphorylation at Ser-202 and
Thr-205 (AT8 epitopes) is the initiating lesion in tau pathology.

The S-entropy mapping is exact: $\St$ measures timing disorder (irregular
membrane potential oscillations = increased $\St$); $\Sk$ measures state
uncertainty (the cell does not know its metabolic state = increased $\Sk$);
$\Se$ measures trajectory loss (a neuron undergoing tangle-mediated death
has lost its developmental trajectory = increased $\Se$). \qed
\end{proof}

\begin{remark}[The Sequence as Partition Cascade]
The Alzheimer cascade is not a collection of independent pathologies.
It is a single partition cascade: the initial breach at the membrane
(L1 partition boundary) propagates energy-deficiency upward through L2--L4
and ultimately into the epigenetic--identity coordinate (L5, $\Se$). Each
step is forced by the previous one because the levels are hierarchically
coupled. This is why no single-target therapy has succeeded: the cascade
initiates below the level at which most drugs intervene.
\end{remark}

%------------------------------------------------------------------------------
\section{Aging as Nuclear Pore Complex Degradation}
\label{sec:aging_npc}
%------------------------------------------------------------------------------

\subsection{NPC Half-Life and Epigenetic Memory}

The nuclear pore complex (NPC) is the sole conduit for macromolecular
exchange between the nucleus and cytoplasm. NPCs are established during
neuronal differentiation and, in post-mitotic neurons, are not replaced for
the lifetime of the cell. Measured NPC half-life in cortical neurons:
$\tau_{\mathrm{NPC}} \approx 20\text{--}30\,\mathrm{years}$.

The NPC is not merely a passive channel; it is a categorical filter that
maintains the partition between the nuclear compartment (where the genome's
chromatin state is stored) and the cytoplasm. Degradation of NPC components
(particularly nucleoporin~Nup98 and Nup88) increases non-selective
permeability, allowing cytoplasmic components to access the nuclear
compartment and disrupting chromatin organisation.

\begin{theorem}[Aging as Monotonic $\Se$ Increase]
\label{thm:aging_se}
In post-mitotic neurons, the evolution entropy coordinate $\Se$ increases
monotonically with chronological age at rate:
\begin{equation}
  \frac{d\Se}{dt} = \frac{\lambda_{\mathrm{NPC}}}{\tau_{\mathrm{NPC}}},
  \label{eq:aging_rate}
\end{equation}
where $\lambda_{\mathrm{NPC}} > 0$ is a proportionality constant encoding
the contribution of one NPC degradation event to $\Se$, and
$\tau_{\mathrm{NPC}} \approx 20\text{--}30\,\mathrm{years}$ is the mean
NPC half-life. This increase is irreversible in the absence of NPC
regeneration (which requires cell division, not available in post-mitotic
neurons).
\end{theorem}

\begin{proof}
$\Se$ measures the normalised uncertainty in the cell's evolutionary
trajectory (Definition~\ref{def:sentropy_space} of
Chapter~\ref{chap:sentropy_space}). The nuclear epigenetic state encodes
$\Se$: a well-compartmentalised nucleus with intact chromatin organisation
has $\Se \approx 0$ (trajectory is determined); a nucleus with degraded
NPC-mediated compartmentalisation has increased cytoplasmic access to
chromatin, leading to aberrant chromatin modifications and loss of gene
expression specificity, increasing $\Se$.

Each NPC degradation event reduces the categorical depth of the nucleus-
cytoplasm partition by one unit. With $N_0 \approx 3{,}000\text{--}5{,}000$
NPCs per neuronal nucleus and a mean degradation rate of $N_0 / \tau_{\mathrm{NPC}}$
events per year, the accumulation rate of partition depth loss is proportional
to $1/\tau_{\mathrm{NPC}}$. Mapping this partition depth loss to the $\Se$
coordinate (which is the normalised partition depth in the evolutionary
trajectory dimension) yields Equation~\eqref{eq:aging_rate}.

Irreversibility follows from the Residue Principle
(Corollary~\ref{cor:residue} of Chapter~\ref{chap:partition_depth}):
each NPC degradation event is a partition boundary erasure that cannot be
undone without energy input exceeding the original construction cost of the
NPC ($\sim 120\,\mathrm{MDa}$ complex, requiring $\sim 10^5\,\mathrm{ATP}$
equivalents for assembly). \qed
\end{proof}

\begin{corollary}[Neurodegeneration as Threshold Crossing]
\label{cor:nd_threshold}
Neurodegenerative disease onset corresponds to the $\Se$ coordinate crossing
a threshold value $\Se^* \approx 0.7$ (empirically corresponding to
$\sim 50\%$ NPC functional integrity loss), at which the nucleus--cytoplasm
partition can no longer maintain chromatin organisation against cytoplasmic
invasion. Disease onset accelerates because NPC loss is autocatalytic:
degraded NPCs cannot efficiently import nuclear proteins needed to maintain
remaining NPCs.
\end{corollary}

\subsection{Connection to Other Neurodegenerative Diseases}

The NPC-degradation mechanism is general across neurodegenerative diseases.
In ALS (amyotrophic lateral sclerosis), TDP-43 and FUS mislocalisation from
nucleus to cytoplasm is a primary pathological feature---explained as loss of
the NPC partition that maintains these RNA-binding proteins in the nucleus.
In Huntington's disease, mutant huntingtin (polyQ-expanded) disrupts Ran-GTPase
activity, the engine of nucleocytoplasmic transport, producing equivalent
$\Se$ elevation through a different molecular mechanism.

%------------------------------------------------------------------------------
\section{The Unified Pathology Theorem}
\label{sec:unified_pathology}
%------------------------------------------------------------------------------

The three disease classes of Sections~\ref{sec:cancer_warburg}--\ref{sec:aging_npc}
share a common formal structure. This structure is not accidental.

\begin{theorem}[Unified Pathology as Depth Reduction]
\label{thm:unified_pathology}
All major chronic diseases reduce hierarchical categorical depth $\Depth$.
Health requires $\Depth = 5$. Chronic disease is defined by $\Depth \leq 4$.
Death occurs when $\Depth < 3$, below the life threshold established in
Chapter~\ref{chap:what_is_life}.
\end{theorem}

\begin{proof}
Each disease class corresponds to a specific locus of depth reduction:
\begin{itemize}
  \item \textbf{Cancer}: L4 decoupling (Theorem~\ref{thm:cancer_depth_collapse})
        reduces $\Depth$ from 5 to 2 in advanced disease.
  \item \textbf{Neurodegeneration}: L1 partition breach (Theorem~\ref{thm:alzheimer_cascade})
        propagates to reduce all levels; eventually $\Depth < 3$.
  \item \textbf{Aging}: $\Se$ increase (Theorem~\ref{thm:aging_se}) reduces
        the effective partition depth of the nucleus--cytoplasm interface,
        equivalent to degradation of L5 (epigenetic--metabolic feedback),
        reducing $\Depth$ from 5 toward 4 over decades.
\end{itemize}

The life threshold $\Depth \geq 3$ is independently established in
Chapter~\ref{chap:what_is_life}: an entity requires at least 3 active levels
of the partition hierarchy to autonomously maintain its boundary, replicate
information, and respond to perturbation. Below $\Depth = 3$, one or more
of these three requirements fails, and the entity dissolves into the
environment in finite time. \qed
\end{proof}

Table~\ref{tab:disease_sentropy} provides a unified classification of major
disease classes in terms of S-entropy deviations.

\begin{table}[htbp]
\centering
\caption{Major disease classes mapped to S-entropy coordinate deviations
         and partition depth reductions.
         $\uparrow$ = elevated above health value;
         $\downarrow$ = reduced below health value.}
\label{tab:disease_sentropy}
\begin{tabularx}{\textwidth}{@{} l c c c c X @{}}
  \toprule
  Disease Class & $\Depth$ & $\Sk$ & $\St$ & $\Se$ & Primary Partition Lesion \\
  \midrule
  Cancer (early) & 4 & $\uparrow$ & $\uparrow$ & neutral & Driver mutation: one L-level compromised \\
  Cancer (Warburg) & 2 & $\uparrow\uparrow$ & $\uparrow\uparrow$ & $\uparrow$ & L3--L4 decoupling by acidification \\
  Alzheimer's (early) & 4 & $\uparrow$ & $\uparrow\uparrow$ & neutral & Membrane partition breach at L1 \\
  Alzheimer's (late) & $<$3 & $\uparrow\uparrow$ & $\uparrow\uparrow$ & $\uparrow\uparrow$ & Full cascade L1$\to$L5 \\
  Parkinson's disease & 3--4 & neutral & $\uparrow\uparrow$ & neutral & Mitochondrial fission excess: L3--L4 \\
  Type 2 diabetes & 4 & $\uparrow\uparrow$ & neutral & neutral & Insulin-partition decoupling at L3 \\
  Atherosclerosis & 4 & neutral & $\uparrow$ & $\uparrow$ & Endothelial partition breach \\
  Normal aging & $\to$4 & neutral & $\uparrow$ & $\uparrow\uparrow$ & NPC degradation: L5 $\Se$ accumulation \\
  \bottomrule
\end{tabularx}
\end{table}

%------------------------------------------------------------------------------
\section{Partition-Based Therapeutics}
\label{sec:partition_therapeutics}
%------------------------------------------------------------------------------

The unified pathology theorem implies a unified therapeutic strategy: rather
than targeting specific molecules downstream of the primary lesion, restore
the partition topology of the affected level.

\begin{definition}[Partition-Targeted Therapy]
\label{def:partition_therapy}
A \emph{partition-targeted therapy} for a disease at depth $\Depth = k$ is
an intervention that:
\begin{enumerate}[label=(\roman*)]
  \item restores $\Depth$ from $k$ to $k+1$ or higher by re-establishing
        the failed partition at level L$(k+1)$;
  \item preserves $Q_{\mathrm{net}}$ (ionic charge homeostasis) in the
        process;
  \item does not create new partition breaches at other levels (which
        would cause side effects).
\end{enumerate}
\end{definition}

\begin{example}[Metformin as L3--L4 Partition Restorer]
\label{ex:metformin}
Metformin (1,1-dimethylbiguanide) is the first-line pharmacological treatment
for type~2 diabetes. Its conventional description is that it ``lowers blood
glucose'', but this is a downstream effect. The partition framework identifies
the primary mechanism: metformin is a mild inhibitor of mitochondrial Complex~I
(NADH:ubiquinone oxidoreductase), the first enzyme of the electron transport
chain (L3).

At clinical concentrations ($\sim 10$--$100\,\mu\mathrm{M}$ in hepatocytes),
metformin reduces Complex~I turnover by $\sim 20\text{--}40\%$. This
partial inhibition, paradoxically, \emph{restores} L3--L4 coupling in
metabolically dysregulated cells:
\begin{enumerate}[label=(\roman*)]
  \item In insulin-resistant hepatocytes, Complex~I is hyperactive,
        generating excess NADH and driving excessive gluconeogenesis via
        the TCA cycle.
  \item Metformin-mediated partial inhibition brings L3 flux into balance
        with L4 (ATP synthase) capacity, re-establishing tight L3--L4 coupling.
  \item The restored coupling increases the efficiency of mitochondrial ATP
        production, reducing the need for compensatory glycolysis.
  \item Downstream glucose output falls as a consequence of restored
        partition depth, not as a primary drug action.
\end{enumerate}
In partition terms: metformin increases $\Depth$ from 4 to 5 in
insulin-resistant hepatocytes by restoring the L3--L4 partition interface,
reducing $\Sk$ (the cell's metabolic state becomes known and coherent).
The blood glucose reduction is a consequence of this $\Depth$ restoration.
\end{example}

\begin{proposition}[Therapeutic Design Principle]
\label{prop:therapy_principle}
For a disease at partition depth $\Depth = k$, the effective therapeutic
window is the set of interventions that:
\begin{equation}
  \max_{\text{interventions}} \left[\Depth(\Scoord_{\text{post-treatment}}) - \Depth(\Scoord_{\text{disease}})\right],
\end{equation}
subject to the constraint that non-target partition levels are perturbed by
less than $\epsilon_{\text{off-target}}$ (the off-target tolerance).
Drugs that modify partition depth at precisely the failed level have the highest
ratio of efficacy to off-target effects.
\end{proposition}

\begin{proof}
The disease state is $\Scoord_D$ with $\Depth(\Scoord_D) = k$.
The health state is $\Scoord_H$ with $\Depth(\Scoord_H) = 5$.
An intervention is effective if it moves $\Scoord$ from $\Scoord_D$ toward
$\Scoord_H$, i.e., if it increases $\Depth$. The maximum increase per
intervention is bounded by $\min(5 - k, \Delta\Depth_{\text{drug}})$, where
$\Delta\Depth_{\text{drug}}$ is the depth change achievable by the drug.

Off-target effects arise when the drug modifies partition depth at levels
other than the target. Since partition levels are hierarchically coupled,
interventions at the correct level propagate upward (toward L5) with
decreasing amplitude, while interventions at the wrong level propagate
sideways (to non-target pathways) without beneficial effect. Targeting the
failed level is therefore both maximally efficacious and minimally off-target
by construction. \qed
\end{proof}

\begin{remark}[Implications for Drug Discovery]
The partition framework reframes drug discovery: the objective is not to find
a molecule that binds a target with high affinity ($K_d$ minimisation), but
to find a molecule that restores $\Depth$ at the correct level with minimal
off-target $\Depth$ modification. These are not the same objective, and
optimising for the former while ignoring the latter is the primary reason
that high-affinity, high-selectivity compounds fail in clinical trials:
they perturb the target partition but do not restore the failed level.
\end{remark}

%------------------------------------------------------------------------------
\section{The Health Attractor and Disease Basins}
\label{sec:health_attractor}
%------------------------------------------------------------------------------

The partition landscape of a healthy cell has a single deep minimum (the
health attractor $\Scoord_H$) and a collection of shallower minima (disease
attractors $\Scoord_{D_i}$) separated by partition ridges.

\begin{proposition}[Basin Depth of Health vs. Disease]
\label{prop:health_basin}
The health attractor $\Scoord_H$ is deeper than any disease attractor
$\Scoord_{D_i}$ in the sense that:
\begin{equation}
  \Depth(\Scoord_H) > \Depth(\Scoord_{D_i}) \quad \text{for all } i,
\end{equation}
because health requires all five levels to be active, while disease activates
at most four levels. The health attractor is therefore more partition-rich
(more distinct functional states are accessible from $\Scoord_H$ than from
any $\Scoord_{D_i}$).
\end{proposition}

\begin{proof}
At $\Scoord_H$, $\Depth = 5$: all five metabolic-genomic coupling levels are
active, and the cell can access $3^5 = 243$ distinct categorical states in the
hierarchical partition. At $\Scoord_{D_i}$ with $\Depth = k < 5$, only
$3^k < 243$ states are accessible. The health state is therefore richer by
any measure of categorical depth. The cell at $\Scoord_H$ has greater
adaptive capacity (more reachable states) than at any disease attractor.
This is the partition-theoretic meaning of robustness or resilience. \qed
\end{proof}

\begin{remark}[Why Cancer Cannot Be Healthier Than Normal]
Proposition~\ref{prop:health_basin} provides a formal proof of an intuitively
obvious statement: cancer cells at $\Depth = 2$ have $3^2 = 9$ accessible
states, compared to $3^5 = 243$ for normal cells. Cancer is not an
alternative healthy state; it is an impoverished state with radically fewer
degrees of freedom. The apparently vigorous proliferation of cancer cells is
the gradient flow toward the L1 minimum, not evidence of increased complexity.
\end{remark}

%==============================================================================
\section*{Chapter Summary}
\label{sec:ch22_summary}
%==============================================================================

\begin{chapterbox}[title=Chapter 22 --- Key Results Established]
\begin{enumerate}[label=\textbf{\arabic*.}]
  \item \textbf{Disease Attractor (Definition~\ref{def:disease_attractor})}:
        Disease is a local minimum of the partition landscape $\Lambda$ with
        $\Depth < 5$. Disease progression is gradient flow from the health
        attractor toward the disease attractor.

  \item \textbf{Phosphate Protonation (Proposition~\ref{prop:phosphate_protonation})}:
        Tumour-microenvironment acidification (pH $6.7$--$7.0$) protonates
        $\sim 39\%$ of DNA phosphate groups, reducing genomic charge by
        $\sim 39\%$ and collapsing L3--L4 electromagnetic coupling.

  \item \textbf{Cancer as Depth Collapse (Theorem~\ref{thm:cancer_depth_collapse})}:
        Cancer stages map to $\Depth = 5 \to 4 \to 2 \to <2$ in a stereotyped
        sequence following the metabolic hierarchy L5 to L1.
        TCGA metabolomics confirms the metabolite signatures of each stage.

  \item \textbf{Alzheimer Cascade (Theorem~\ref{thm:alzheimer_cascade})}:
        A$\beta$ oligomers initiate a seven-step partition cascade from L1
        membrane breach through ATP depletion, Ca$^{2+}$ overload, tau
        phosphorylation, and synaptic failure. The cascade maps onto
        monotonically increasing $\Sk$, $\St$, and $\Se$.

  \item \textbf{Aging as $\Se$ Accumulation (Theorem~\ref{thm:aging_se})}:
        In post-mitotic neurons, $\Se$ increases at rate $\lambda_{\mathrm{NPC}}/\tau_{\mathrm{NPC}}$
        due to irreversible NPC degradation over a half-life of $20$--$30$ years.

  \item \textbf{Unified Pathology Theorem (Theorem~\ref{thm:unified_pathology})}:
        All major diseases reduce $\Depth$. Death occurs at $\Depth < 3$,
        consistent with the life threshold of Chapter~\ref{chap:what_is_life}.

  \item \textbf{Partition-Targeted Therapy (Definition~\ref{def:partition_therapy}
        and Proposition~\ref{prop:therapy_principle})}:
        Effective therapy restores $\Depth$ at the failed level.
        Metformin exemplifies this: it restores L3--L4 coupling depth in
        insulin-resistant cells, with blood-glucose reduction as a downstream
        consequence.

  \item \textbf{Health Basin Depth (Proposition~\ref{prop:health_basin})}:
        $\Depth = 5$ gives $3^5 = 243$ accessible categorical states versus
        $3^2 = 9$ at $\Depth = 2$ (advanced cancer). Health is not merely
        the absence of disease but the most partition-rich attractor in the
        landscape.
\end{enumerate}
\end{chapterbox}
